%\documentclass[letterpaper,12pt]{article}
\documentclass[preview]{standalone}

\usepackage{tabularx}
\usepackage{amsmath}
\usepackage{graphicx}
\usepackage[margin=1in,letterpaper]{geometry}
\usepackage{cite}
\usepackage[final]{hyperref}
\usepackage{datatool}
\usepackage{caption}
\usepackage{float}
\usepackage{gensymb}

\hypersetup{
        colorlinks=true,
        linkcolor=blue,
        citecolor=blue,
        filecolor=magenta,
        urlcolor=blue
}
\begin{document}

\setcounter{table}{2}
\begin{table}[h!]
\centering
\label{tab:fits}
\renewcommand{\arraystretch}{1.2}
\begin{tabular}{|c|c|c|c|c|c|}
\hline
R $\pm 1$ $\Omega$ & $A$ | V & $B$ | V & $\lambda$ & $\omega$ | rad/s & $\phi$ | $^\circ$ \\
\hline
839 &
$2.625 \pm 0.004$ &
$-2384 \pm 16$ &
$1595 \pm 10$ &
-- & -- \\

10   & 
$4.633 \pm 0.005$ &
-- &
$179.5 \pm 0.3$ &
$4455.8 \pm 0.3$ &
$0.917 \pm 0.001$ \\

500   & 
$2.949 \pm 0.005$ &
-- &
$2459 \pm 11$ &
$3541 \pm 15$ &
$0.401 \pm 0.004$ \\

1000   &
$2.645 \pm 0.004$ &
$-1930 \pm 14$ &
$1472 \pm 7$ &
-- & -- \\

3500   &
$-0.044 \pm 0.006$ &
$2.82 \pm 0.01$ &
wtf &
-- \\
\hline
\end{tabular}
\caption{Paramètres ajustés des différents régimes d'amortissement}
\end{table}

\end{document}

